\chapter{The First Appearance of Colbert}

\lettrine{T}{he} whole night was passed in anguish, common to the dying man and to the king: the dying man expected his deliverance, the king awaited his liberty. Louis did not go to bed. An hour after leaving the chamber of the cardinal, he learned that the dying man, recovering a little strength, had insisted upon being dressed, adorned and painted, and seeing the ambassadors. Like Augustus, he no doubt considered the world a great stage, and was desirous of playing out the last act of the comedy. Anne of Austria reappeared no more in the cardinal's apartments; she had nothing more to do there. Propriety was the pretext for her absence. On his part, the cardinal did not ask for her: the advice the queen had giver her son rankled in his heart.

Towards midnight, while still painted, Mazarin's mortal agony came on. He had revised his will, and as this will was the exact expression of his wishes, and as he feared that some interested influence might take advantage of his weakness to make him change something in it, he had given orders to Colbert, who walked up and down the corridor which led to the cardinal's bed-chamber, like the most vigilant of sentinels. The king, shut up in his own apartment, dispatched his nurse every hour to Mazarin's chamber, with orders to bring him back an exact bulletin of the cardinal's state. After having heard that Mazarin was dressed, painted, and had seen the ambassadors, Louis herd that the prayers for the dying were being read for the cardinal. At one o'clock in the morning, Guénaud had administered the last remedy. This was a relic of the old customs of that fencing time, which was about to disappear to give place to another time, to believe that death could be kept off by some good secret thrust. Mazarin, after having taken the remedy, respired freely for nearly ten minutes. He immediately gave orders that the news should be spread everywhere of a fortunate crisis. The king, on learning this, felt as if a cold sweat were passing over his brow;—he had had a glimpse of the light of liberty; slavery appeared to him more dark and less acceptable than ever. But the bulletin which followed entirely changed the face of things. Mazarin could no longer breathe at all, and could scarcely follow the prayers which the cure of Saint-Nicholas-des-Champs recited near him. The king resumed his agitated walk about his chamber, and consulted, as he walked, several papers drawn from a casket of which he alone had the key. A third time the nurse returned. M.~de Mazarin had just uttered a joke, and had ordered his \workname{Flora,} by Titian, to be revarnished. At length, towards two o'clock in the morning, the king could no longer resist his weariness: he had not slept for twenty-four hours. Sleep, so powerful at his age, overcame him for about an hour. But he did not go to bed for that hour; he slept in a fauteuil. About four o'clock his nurse awoke him by entering the room.

<Well?> asked the king.

<Well, my dear sire,> said the nurse, clasping her hands with an air of commiseration. <Well; he is dead!>

The king arose at a bound, as if a steel spring had been applied to his legs. <Dead!> cried he.

<Alas! yes.>

<Is it quite certain?>

<Yes.>

<Official?>

<Yes.>

<Has the news been made public?>

<Not yet.>

<Who told you, then, that the cardinal was dead?>

<M.~Colbert.>

<M.~Colbert?>

<Yes.>

<And he was sure of what he said?>

<He came out of the chamber, and had held a glass for some minutes before the cardinal's lips.>

<Ah!> said the king. <And what is become of M.~ Colbert?>

<He has just left his eminence's chamber.>

<Where is he?>

<He followed me.>

<So that he is\longdash>

<Sire, waiting at your door, till it shall be your good pleasure to receive him.>

Louis ran to the door, opened it himself, and perceived Colbert standing waiting in the passage. The king started at sight of this statue, all clothed in black. Colbert, bowing with profound respect, advanced two steps towards his majesty. Louis re-entered his chamber, making Colbert a sign to follow. Colbert entered; Louis dismissed the nurse, who closed the door as she went out. Colbert remained modestly standing near that door.

<What do you come to announce to me, monsieur?> said Louis, very much troubled at being thus surprised in his private thoughts, which he could not completely conceal.

<That monsieur le cardinal has just expired, sire; and that I bring your majesty his last adieu.>

The king remained pensive for a minute; and during that minute he looked attentively at Colbert;—it was evident that the cardinal's last words were in his mind. <Are you, then, M.~Colbert?> asked he.

<Yes, sire.>

<His faithful servant, as his eminence himself told me?>

<Yes, sire.>

<The depositary of many of his secrets?>

<Of all of them.>

<The friends and servants of his eminence will be dear to me, monsieur, and I shall take care that you are well placed in my employment.>

Colbert bowed.

<You are a financier, monsieur, I believe?>

<Yes, sire.>

<And did monsieur le cardinal employ you in his stewardship?>

<I had that honour, sire.>

<You never did anything personally for my household, I believe?>

<Pardon me, sire, it was I who had the honour of giving monsieur le cardinal the idea of an economy which puts three hundred thousand francs a year into your majesty's coffers.>

<What economy was that, monsieur?> asked \regnal{Louis XIV}.

<Your majesty knows that the hundred Swiss have silver lace on each side of their ribbons?>

<Doubtless.>

<Well, sire, it was I who proposed that imitation silver lace should be placed upon these ribbons; it could not be detected, and a hundred thousand crowns serve to feed a regiment during six months; and is the price of ten thousand good muskets or the value of a vessel of ten guns, ready for sea.>

<That is true,> said \regnal{Louis XIV}, considering more attentively, <and, \swear{ma foi!} that was a well placed economy; besides, it was ridiculous for soldiers to wear the same lace as noblemen.>

<I am happy to be approved of by your majesty.>

<Is that the only appointment you held about the cardinal?> asked the king.

<It was I who was appointed to examine the accounts of the superintendent, sire.>

<Ah!> said Louis, who was about to dismiss Colbert, but whom that word stopped; <ah! it was you whom his eminence had charged to control M.~ Fouquet, was it? And the result of that examination?>

<Is that there is a deficit, sire; but if your majesty will permit me\longdash>

<Speak, M.~Colbert.>

<I ought to give your majesty some explanations.>

<Not at all, monsieur, it is you who have controlled these accounts; give me the result.>

<That is very easily done, sire: emptiness everywhere, money nowhere.>

<Beware, monsieur; you are roughly attacking the administration of M.~ Fouquet, who, nevertheless, I have heard say, is an able man.>

Colbert coloured, and then became pale, for he felt that from that minute he entered upon a struggle with a man whose power almost equalled the sway of him who had just died. <Yes, sire, a very able man,> repeated Colbert, bowing.

<But if M.~Fouquet is an able man, and, in spite of that ability, if money be wanting, whose fault is it?>

<I do not accuse, sire, I verify.>

<That is well; make out your accounts, and present them to me. There is a deficit, you say? A deficit may be temporary; credit returns and funds are restored.>

<No, sire.>

<Upon this year, perhaps, I understand that; but upon next year?>

<Next year is eaten as bare as the current year.>

<But the year after, then?>

<Will be just like next year.>

<What do you tell me, Monsieur Colbert?>

<I say there are four years engaged beforehand.>

<They must have a loan, then.>

<They must have three, sire.>

<I will create offices to make them resign, and the salary of the posts shall be paid into the treasury.>

<Impossible, sire, for there have already been creations upon creations of offices, the provisions of which are given in blank, so that the purchasers enjoy them without filling them. That is why your majesty cannot make them resign. Further, upon each agreement M.~Fouquet has made an abatement of a third, so that the people have been plundered, without your majesty profiting by it.>

The king started. <Explain me that, M.~Colbert,> he said.

<Let your majesty set down clearly your thought, and tell me what you wish me to explain.>

<You are right, clearness is what you wish, is it not?>

<Yes, sire, clearness. God is God above all things, because He made light.>

<Well, for example,> resumed \regnal{Louis XIV}, <if to-day, the cardinal being dead, and I being king, suppose I wanted money?>

<Your majesty would not have any.>

<Oh! that is strange, monsieur! How! my superintendent would not find me any money?>

Colbert shook his large head.

<How is that?> said the king; <is the income of the state so much in debt that there is no longer any revenue?>

<Yes, sire.>

The king frowned and said, <If it be so, I will get together the ordonnances to obtain a discharge from the holders, a liquidation at a cheap rate.>

<Impossible, for the ordonnances have been converted into bills, which bills, for the convenience of return and facility of transaction, are divided into so many parts that the originals can no longer be recognized.>

Louis, very much agitated, walked about, still frowning. <But, if this is as you say, Monsieur Colbert,> said he, stopping all at once, <I shall be ruined before I begin to reign.>

<You are, in fact, sire,> said the impassible caster-up of figures.

<Well, but yet, monsieur, the money is somewhere?>

<Yes, sire, and even as a beginning, I bring your majesty a note of funds which M.~le Cardinal Mazarin was not willing to set down in his testament, neither in any act whatever, but which he confided to me.>

<To you?>

<Yes, sire, with an injunction to remit it to your majesty.>

<What! besides the forty millions of the testament?>

<Yes, sire.>

<M.~de Mazarin had still other funds?>

Colbert bowed.

<Why, that man was a gulf!> murmured the king. <M.~de Mazarin on one side, M.~Fouquet on the other,—more than a hundred millions perhaps between them! No wonder my coffers should be empty!> Colbert waited without stirring.

<And is the sum you bring me worth the trouble?> asked the king.

<Yes, sire, it is a round sum.>

<Amounting to how much?>

<To thirteen millions of livres, sire.>

<Thirteen millions!> cried Louis, trembling with joy; <do you say thirteen millions, Monsieur Colbert?>

<I said thirteen millions, yes, your majesty.>

<Of which everybody is ignorant?>

<Of which everybody is ignorant.>

<Which are in your hands?>

<In my hands, yes, sire.>

<And which I can have?>

<Within two hours, sire.>

<But where are they, then?>

<In the cellar of a house which the cardinal possessed in the city, and which he was so kind as to leave me by a particular clause of his will.>

<You are acquainted with the cardinal's will, then?>

<I have a duplicate of it, signed by his hand.>

<A duplicate?>

<Yes, sire, and here it is.> Colbert drew the deed quietly from his pocket, and showed it to the king. The king read the article relative to the donation of the house.

<But,> said he, <there is no question here but of the house; there is nothing said of the money.>

<Your pardon, sire, it is in my conscience.>

<And Monsieur Mazarin has intrusted it to you?>

<Why not, sire?>

<He! a man mistrustful of everybody?>

<He was not so of me, sire, as your majesty may perceive.>

Louis fixed his eyes with admiration upon that vulgar but expressive face. <You are an honest man, M.~Colbert,> said the king.

<That is not a virtue, it is a duty,> replied Colbert, coolly.

<But,> added Louis, <does not the money belong to the family?>

<If this money belonged to the family it would be disposed of in the testament, as the rest of the fortune is. If this money belonged to the family, I, who drew up the deed of donation in favour of your majesty, should have added the sum of thirteen millions to that of forty millions which was offered to you.>

<How!> exclaimed \regnal{Louis XIV}, <was it you who drew up the deed of donation?>

<Yes, sire.>

<And yet the cardinal was attached to you?> added the king, ingenuously.

<I had assured his eminence you would by no means accept the gift,> said Colbert, in that same quiet manner we have described, and which, even in the common habits of life, had something solemn in it.

Louis passed his hand over his brow: <Oh! how young I am,> murmured he, <to have command of men.>

Colbert waited the end of this monologue. He saw Louis raise his head. <At what hour shall I send the money to your majesty?> asked he.

<To-night, at eleven o'clock; I desire that no one may know that I possess this money.>

Colbert made no more reply than if the thing had not been said to him.

<Is the amount in ingots, or coined gold?>

<In coined gold, sire.>

<That is well.>

<Where shall I send it?>

<To the Louvre. Thank you, M.~Colbert.>

Colbert bowed and retired. <Thirteen millions!> exclaimed Louis, as soon as he was alone. <This must be a dream!> Then he allowed his head to sink between his hands, as if he were really asleep. But, at the end of a moment, he arose, and opening the window violently, he bathed his burning brow in the keen morning air, which brought to his senses the scent of the trees, and the perfume of the flowers. A splendid dawn was gilding the horizon, and the first rays of the sun bathed in flame the young king's brow. <This is the dawn of my reign,> murmured \regnal{Louis XIV}. <It's a presage sent by the Almighty.> 